\section{FODA de algoritmos de planificación de CPU}

\subsection{FCFS -- First-Come, First-Served}

\begin{center}
\begin{tabular}{|p{0.21\textwidth}|p{0.67\textwidth}|}
\hline
\textbf{FODA} & \textbf{Descripción} \\ \hline

\textbf{Fortalezas} &
Es un algoritmo sencillo de implementar y comprender. Las tareas se ejecutan en el mismo orden en que llegan, por lo que su funcionamiento es predecible y no requiere cálculos adicionales para decidir cuál proceso ejecutar. Además, evita la inanición, ya que todo proceso eventualmente obtiene acceso al CPU. \\ \hline

\textbf{Oportunidades} &
Puede utilizarse en sistemas donde la simplicidad y la facilidad de administración sean más importantes que el tiempo de respuesta. También sirve como base para comprender algoritmos de planificación más avanzados. \\ \hline

\textbf{Debilidades} &
Puede producir tiempos de espera elevados cuando una tarea larga llega antes que varias tareas cortas. Este comportamiento se conoce como \textit{efecto convoy}. No es adecuado para sistemas interactivos donde se necesita una respuesta rápida. \\ \hline

\textbf{Amenazas} &
En sistemas con muchas tareas o con procesos de duración muy diferente, el tiempo de respuesta puede aumentar considerablemente. Esto puede provocar una experiencia poco fluida para el usuario. \\ \hline

\end{tabular}
\end{center}

\subsection{SJF -- Shortest-Job-First}

\begin{center}
\begin{tabular}{|p{0.21\textwidth}|p{0.67\textwidth}|}
\hline
\textbf{FODA} & \textbf{Descripción} \\ \hline

\textbf{Fortalezas} &
Reduce considerablemente el tiempo promedio de espera cuando se conocen las duraciones de las tareas. En condiciones ideales, proporciona un promedio de espera menor que FCFS. \\ \hline

\textbf{Oportunidades} &
Es especialmente útil en sistemas donde las tareas tienen duraciones relativamente previsibles. También constituye la base conceptual de algoritmos como SRTF, que incorpora expulsión de procesos. \\ \hline

\textbf{Debilidades} &
Necesita conocer o estimar previamente cuánto tiempo requerirá cada proceso. Las tareas largas pueden permanecer esperando durante mucho tiempo si continuamente llegan tareas más cortas. \\ \hline

\textbf{Amenazas} &
Puede producir \textit{inanición} (\textit{starvation}) de procesos largos en sistemas donde llegan continuamente tareas cortas. Además, una estimación incorrecta de la duración de los procesos puede reducir sus beneficios. \\ \hline

\end{tabular}
\end{center}

\subsection{SRTF -- Shortest-Remaining-Time-First}

\begin{center}
\begin{tabular}{|p{0.21\textwidth}|p{0.67\textwidth}|}
\hline
\textbf{FODA} & \textbf{Descripción} \\ \hline

\textbf{Fortalezas} &
Es una versión expulsiva de SJF que selecciona el proceso con menor tiempo de ejecución restante. Puede ofrecer tiempos de espera y respuesta muy bajos para tareas cortas, especialmente cuando las duraciones de los procesos pueden estimarse correctamente. \\ \hline

\textbf{Oportunidades} &
Puede ser útil en sistemas donde existen numerosas tareas cortas y se busca mejorar el tiempo de respuesta. Su capacidad de interrumpir un proceso permite adaptarse dinámicamente a la llegada de nuevas tareas. \\ \hline

\textbf{Debilidades} &
Requiere un mayor control del sistema, ya que puede producir múltiples cambios de contexto. Además, necesita conocer o estimar el tiempo restante de ejecución de los procesos. \\ \hline

\textbf{Amenazas} &
Los procesos largos pueden sufrir inanición si llegan continuamente procesos con menor tiempo restante. Un exceso de cambios de contexto también puede aumentar la sobrecarga del sistema y disminuir su eficiencia. \\ \hline

\end{tabular}
\end{center}

\subsection{RR -- Round-Robin}

\begin{center}
\begin{tabular}{|p{0.21\textwidth}|p{0.67\textwidth}|}
\hline
\textbf{FODA} & \textbf{Descripción} \\ \hline

\textbf{Fortalezas} &
Es un algoritmo expulsivo diseñado especialmente para sistemas interactivos. Cada proceso recibe un intervalo de tiempo denominado \textit{quantum}, lo que permite repartir el uso del CPU de manera equitativa y garantizar una buena capacidad de respuesta. \\ \hline

\textbf{Oportunidades} &
Es apropiado para sistemas multiusuario, entornos interactivos y situaciones donde varias tareas deben recibir atención periódicamente. Puede ajustarse mediante la elección adecuada del tamaño del \textit{quantum}. \\ \hline

\textbf{Debilidades} &
Su eficiencia depende del tamaño del \textit{quantum}. Si es demasiado pequeño, aumenta el número de cambios de contexto; si es demasiado grande, el algoritmo se comporta de manera similar a FCFS y puede empeorar el tiempo de respuesta. \\ \hline

\textbf{Amenazas} &
Una elección inadecuada del \textit{quantum} puede generar sobrecarga o tiempos de respuesta poco satisfactorios. En cargas de trabajo con procesos muy cortos o muy largos, puede no alcanzar el mejor tiempo promedio de espera. \\ \hline

\end{tabular}
\end{center}
